\documentclass{amsart}
% For dealing with references we use the comment environment
\usepackage{verbatim}
\newenvironment{reference}{\comment}{\endcomment}
%\newenvironment{reference}{}{}
\newenvironment{slogan}{\comment}{\endcomment}
\newenvironment{history}{\comment}{\endcomment}

% For commutative diagrams we use Xy-pic
\usepackage[all]{xy}

% We use 2cell for 2-commutative diagrams.
\xyoption{2cell}
\UseAllTwocells

% We use multicol for the list of chapters between chapters
\usepackage{multicol}

% This is generally recommended for better output
\usepackage{lmodern}
\usepackage[T1]{fontenc}

% For cross-file-references

% Package for hypertext links:
\usepackage{hyperref}

% For any local file, say "hello.tex" you want to link to please

% Theorem environments.
%
\theoremstyle{plain}
\newtheorem{theorem}[subsection]{Theorem}
\newtheorem{proposition}[subsection]{Proposition}
\newtheorem{lemma}[subsection]{Lemma}

\theoremstyle{definition}
\newtheorem{definition}[subsection]{Definition}
\newtheorem{example}[subsection]{Example}
\newtheorem{exercise}[subsection]{Exercise}
\newtheorem{situation}[subsection]{Situation}

\theoremstyle{remark}
\newtheorem{remark}[subsection]{Remark}
\newtheorem{remarks}[subsection]{Remarks}

\numberwithin{equation}{subsection}

% Macros
%
\def\lim{\mathop{\mathrm{lim}}\nolimits}
\def\colim{\mathop{\mathrm{colim}}\nolimits}
\def\Spec{\mathop{\mathrm{Spec}}}
\def\Hom{\mathop{\mathrm{Hom}}\nolimits}
\def\Ext{\mathop{\mathrm{Ext}}\nolimits}
\def\SheafHom{\mathop{\mathcal{H}\!\mathit{om}}\nolimits}
\def\SheafExt{\mathop{\mathcal{E}\!\mathit{xt}}\nolimits}
\def\Sch{\mathit{Sch}}
\def\Mor{\mathop{\mathrm{Mor}}\nolimits}
\def\Ob{\mathop{\mathrm{Ob}}\nolimits}
\def\Sh{\mathop{\mathit{Sh}}\nolimits}
\def\NL{\mathop{N\!L}\nolimits}
\def\CH{\mathop{\mathrm{CH}}\nolimits}
\def\proetale{{pro\text{-}\acute{e}tale}}
\def\etale{{\acute{e}tale}}
\def\QCoh{\mathit{QCoh}}
\def\Ker{\mathop{\mathrm{Ker}}}
\def\Im{\mathop{\mathrm{Im}}}
\def\Coker{\mathop{\mathrm{Coker}}}
\def\Coim{\mathop{\mathrm{Coim}}}

% Boxtimes
%
\DeclareMathSymbol{\boxtimes}{\mathbin}{AMSa}{"02}

%
% Macros for moduli stacks/spaces
%
\def\QCohstack{\mathcal{QC}\!\mathit{oh}}
\def\Cohstack{\mathcal{C}\!\mathit{oh}}
\def\Spacesstack{\mathcal{S}\!\mathit{paces}}
\def\Quotfunctor{\mathrm{Quot}}
\def\Hilbfunctor{\mathrm{Hilb}}
\def\Curvesstack{\mathcal{C}\!\mathit{urves}}
\def\Polarizedstack{\mathcal{P}\!\mathit{olarized}}
\def\Complexesstack{\mathcal{C}\!\mathit{omplexes}}
% \Pic is the operator that assigns to X its picard group, usage \Pic(X)
% \Picardstack_{X/B} denotes the Picard stack of X over B
% \Picardfunctor_{X/B} denotes the Picard functor of X over B
\def\Pic{\mathop{\mathrm{Pic}}\nolimits}
\def\Picardstack{\mathcal{P}\!\mathit{ic}}
\def\Picardfunctor{\mathrm{Pic}}
\def\Deformationcategory{\mathcal{D}\!\mathit{ef}}

\setlength{\emergencystretch}{3em}
\begin{document}
\title[Stacks Project, Tag 0BM9]{321 --- Finite etale covers of normal punctured spectra are equivalent to finite normal algebras}
\maketitle
\noindent Copyright (C) 2005--2025 Johan de Jong. Stacks Project authors. Source: \href{https://stacks.math.columbia.edu/tag/0BM9}{Tag 0BM9}.

\noindent Editorial difficulty note (not author text). Difficulty H2: The global sections of a finite etale cover must be proved finite over the local base, despite the punctured spectrum being nonaffine. The proof identifies these sections with an integral closure by localizing along the puncture and using reflexivity and depth-two Hartogs extension, thereby constructing the inverse equivalence..

\noindent Editorial provenance note (not author text). History: excerpt from revision \texttt{a0444\allowbreak 6e57e\allowbreak c1fbc\allowbreak 252a8\allowbreak 71afc\allowbreak ec775\allowbreak 2fb28\allowbreak 07b14}, pione.tex, lines 4792--4857. Reference presentation was adapted; original-author context and core support blocks were retained or added. The mathematical wording is unchanged. Licensed under GNU FDL 1.2 or later; no invariant sections or cover texts. See ../licenses/COPYING. Context blocks reproduce author definitions and supporting results; they are not additional counted problems.

\begin{lemma}
\label{lemma-reformulate-purity-normal}
Let $(A, \mathfrak m)$ be a Noetherian local ring. Set $X = \Spec(A)$
and let $U = X \setminus \{\mathfrak m\}$. Assume $A$ is normal
of dimension $\geq 2$. The functor
$$
\textit{F\'Et}_U \longrightarrow
\left\{
\begin{matrix}
\text{finite normal }A\text{-algebras }B\text{ such} \\
\text{that }\Spec(B) \to X\text{ is \'etale over }U
\end{matrix}
\right\},
\quad
V \longmapsto \Gamma(V, \mathcal{O}_V)
$$
is an equivalence. Moreover, $V = Y \times_X U$ for some $Y \to X$
finite \'etale if and only if $B = \Gamma(V, \mathcal{O}_V)$
is finite \'etale over $A$.
\end{lemma}

\begin{proof}
Observe that $\text{depth}(A) \geq 2$ because $A$ is normal
(Serre's criterion for normality, Algebra, Lemma
\href{https://stacks.math.columbia.edu/tag/031S}{lemma criterion normal (Tag 031S)}).
Thus the final statement follows from Lemma \href{https://stacks.math.columbia.edu/tag/0BLK}{lemma reformulate purity (Tag 0BLK)}.
Given $\pi : V \to U$ finite \'etale, set $B = \Gamma(V, \mathcal{O}_V)$.
If we can show that $B$ is normal and finite over $A$, then
we obtain the displayed functor. Since there is an obvious
quasi-inverse functor, this is also all that we have to show.

\medskip\noindent
Since $A$ is normal, the scheme $V$ is normal
(Descent, Lemma \href{https://stacks.math.columbia.edu/tag/034F}{lemma normal local smooth (Tag 034F)}).
Hence $V$ is a finite disjoint union of integral schemes
(Properties, Lemma \href{https://stacks.math.columbia.edu/tag/033M}{lemma normal Noetherian (Tag 033M)}).
Thus we may assume $V$ is integral.
In this case the function field $L$ of $V$
(Morphisms, Section \href{https://stacks.math.columbia.edu/tag/01RR}{section rational maps (Tag 01RR)})
is a finite separable extension of the fraction field of $A$
(because we get it by looking at the generic fibre
of $V \to U$ and using Morphisms, Lemma
\href{https://stacks.math.columbia.edu/tag/02GL}{lemma etale over field (Tag 02GL)}).
By Algebra, Lemma
\href{https://stacks.math.columbia.edu/tag/032L}{lemma Noetherian normal domain finite separable extension (Tag 032L)}
the integral closure $B' \subset L$ of $A$ in $L$ is finite over $A$.
By More on Algebra, Lemma \href{https://stacks.math.columbia.edu/tag/0BM4}{lemma integral closure reflexive (Tag 0BM4)}
we see that $B'$ is a reflexive $A$-module, which in turn implies
that $\text{depth}_A(B') \geq 2$ by
More on Algebra, Lemma \href{https://stacks.math.columbia.edu/tag/0AVB}{lemma reflexive over normal (Tag 0AVB)}.

\medskip\noindent
Let $f \in \mathfrak m$. Then $B_f = \Gamma(V \times_U D(f), \mathcal{O}_V)$
(Properties, Lemma \href{https://stacks.math.columbia.edu/tag/01P7}{lemma invert f sections (Tag 01P7)}).
Hence $B'_f = B_f$ because $B_f$ is normal (see above),
finite over $A_f$ with fraction field $L$.
It follows that $V = \Spec(B') \times_X U$.
Then we conclude that $B = B'$ from
Lemma \href{https://stacks.math.columbia.edu/tag/0BM8}{lemma sections over punctured spec (Tag 0BM8)}
applied to $\Spec(B') \to X$.
This lemma applies because the localizations $B'_{\mathfrak m'}$
of $B'$ at maximal ideals $\mathfrak m' \subset B'$ lying over
$\mathfrak m$ have depth $\geq 2$ by
Algebra, Lemma \href{https://stacks.math.columbia.edu/tag/0AUK}{lemma depth goes down finite (Tag 0AUK)}
and the remark on depth in the preceding paragraph.
\end{proof}
\end{document}
